
\section{Adversarial Generative Modeling}

Sections~5.1--5.3 developed generative modeling through explicit densities, latent-variable models, and variational inference.
Adversarial methods take a different route.
They define a sampler
\[
x = G_\theta(z), \qquad z\sim p(z),
\]
and train it not by maximizing a tractable likelihood, but by forcing its samples to become hard to distinguish from real data.
This makes adversarial modeling conceptually important: it shows that one can pursue distribution matching even when likelihood evaluation is unavailable or inconvenient.

\medskip

\noindent
\textbf{Section role.}
This section develops the classical GAN objective in textbook form.
It defines the minimax game, proves the formula for the optimal discriminator, derives the generator's induced divergence objective, explains the Jensen--Shannon interpretation, and then separates two important extensions: conditional GANs and the Wasserstein viewpoint.
It closes with an explicit account of instability mechanisms, because the elegance of the theory should not obscure the practical difficulty of the method.

\subsection*{The classical minimax objective}

Let $p_{\mathrm{data}}$ denote the data distribution on $\mathcal X$.
Let $p(z)$ be a simple latent prior, typically a standard Gaussian or a uniform distribution, and let $G_\theta: \mathcal Z\to \mathcal X$ be a generator network.
The generator induces a model distribution $p_G$ on $\mathcal X$ through the sampling rule
\[
z\sim p(z), \qquad x = G_\theta(z).
\]
A discriminator $D_\phi: \mathcal X\to (0,1)$ is then trained to output the probability that an input came from real data rather than the generator.

The classical GAN objective is
\[
\min_\theta \max_\phi V(D_\phi,G_\theta)
=
\min_\theta \max_\phi
\left[
\mathbb E_{x\sim p_{\mathrm{data}}}[\log D_\phi(x)]
+
\mathbb E_{x\sim p_G}[\log(1-D_\phi(x))]
\right].
\]
Equivalently, writing the generator expectation through the latent prior,
\[
V(D_\phi,G_\theta)
=
\mathbb E_{x\sim p_{\mathrm{data}}}[\log D_\phi(x)]
+
\mathbb E_{z\sim p(z)}[\log(1-D_\phi(G_\theta(z)))].
\]

The discriminator tries to maximize classification accuracy between real and generated samples.
The generator tries to minimize the same criterion, thereby pushing generated samples toward the real data manifold.
This is a two-player zero-sum game rather than an ordinary loss minimization problem.

\begin{figure}[ht]
\centering
\setlength{\fboxsep}{7pt}
\renewcommand{\arraystretch}{1.2}
\begin{tabular}{c}
\fbox{\parbox{0.84\textwidth}{\centering
Real data $x\sim p_{\mathrm{data}}$ and latent noise $z\sim p(z)$ feed two branches.
}}\\[0.5em]
$\Downarrow$\\[0.5em]
\begin{tabular}{cc}
\fbox{\parbox{0.37\textwidth}{\centering
Generator branch\\[0.3em]
$z\mapsto G_\theta(z)$ produces fake sample $\tilde x$
}}
&
\fbox{\parbox{0.37\textwidth}{\centering
Data branch\\[0.3em]
real sample $x\sim p_{\mathrm{data}}$
}}
\end{tabular}\\[0.8em]
$\Downarrow$ both are shown to the discriminator $D_\phi$\\[0.5em]
\fbox{\parbox{0.84\textwidth}{\centering
Discriminator outputs $D_\phi(x)$ or $D_\phi(\tilde x)$ and sends gradients back.
The discriminator maximizes real/fake discrimination; the generator updates to make $\tilde x$ look real.
}}
\end{tabular}
\caption{Generator--discriminator interaction in a classical GAN. The discriminator receives both real and generated samples, while the generator is updated only through the discriminator's feedback.}
\label{fig:gan-interaction}
\end{figure}

\subsection*{Optimal discriminator for a fixed generator}

For fixed generator distribution $p_G$, the discriminator optimization can be solved pointwise.
This is the first key theoretical fact.

\begin{proposition}[Optimal discriminator]
Fix a generator distribution $p_G$.
For each $x$ such that $p_{\mathrm{data}}(x)+p_G(x)>0$, the value function is maximized by
\[
D^*(x)=\frac{p_{\mathrm{data}}(x)}{p_{\mathrm{data}}(x)+p_G(x)}.
\]
At points where both densities vanish, the value is irrelevant to the objective.
\end{proposition}

\begin{proof}
For fixed $p_G$, the objective can be written as
\[
V(D,G)
=
\int \left[p_{\mathrm{data}}(x)\log D(x)+p_G(x)\log(1-D(x))\right]dx.
\]
The integrand depends on $D(x)$ pointwise.
For a fixed $x$, define
\[
f_x(d)=a\log d+b\log(1-d),
\]
where $a=p_{\mathrm{data}}(x)$ and $b=p_G(x)$.
Then
\[
f_x'(d)=\frac{a}{d}-\frac{b}{1-d}.
\]
Setting the derivative to zero gives
\[
\frac{a}{d}=\frac{b}{1-d}
\quad \Longrightarrow \quad
 d = \frac{a}{a+b}.
\]
Since
\[
f_x''(d)=-\frac{a}{d^2}-\frac{b}{(1-d)^2}<0,
\]
this critical point is the maximizer.
Substituting back yields the stated formula.
\end{proof}

This proposition makes the discriminator statistically interpretable.
At optimum it returns the posterior probability that a point came from the data rather than the generator under a balanced mixture model.
So the discriminator is not arbitrary: it estimates relative density evidence.

\subsection*{The generator's induced divergence objective}

Once the optimal discriminator is substituted back into the game, the minimax objective becomes a pure distribution-comparison criterion.

\begin{proposition}[Value at the optimal discriminator]
For fixed generator distribution $p_G$ and optimal discriminator $D^*$,
\[
V(D^*,G)
=
\int p_{\mathrm{data}}(x)
\log \frac{p_{\mathrm{data}}(x)}{p_{\mathrm{data}}(x)+p_G(x)}\,dx
+
\int p_G(x)
\log \frac{p_G(x)}{p_{\mathrm{data}}(x)+p_G(x)}\,dx.
\]
Equivalently, if
\[
m(x)=\frac{1}{2}\bigl(p_{\mathrm{data}}(x)+p_G(x)\bigr),
\]
then
\[
V(D^*,G) = -\log 4 + 2\,\mathrm{JS}(p_{\mathrm{data}}\,\|\,p_G),
\]
where $\mathrm{JS}$ denotes the Jensen--Shannon divergence.
\end{proposition}

\begin{proof}[Proof sketch]
Insert the formula for $D^*$ into the value function:
\[
V(D^*,G)
=
\mathbb E_{p_{\mathrm{data}}}
\left[\log \frac{p_{\mathrm{data}}}{p_{\mathrm{data}}+p_G}\right]
+
\mathbb E_{p_G}
\left[\log \frac{p_G}{p_{\mathrm{data}}+p_G}\right].
\]
Now use
\[
p_{\mathrm{data}}+p_G = 2m
\]
and split the logarithms:
\[
\log \frac{p_{\mathrm{data}}}{p_{\mathrm{data}}+p_G}
=
\log \frac{p_{\mathrm{data}}}{m} - \log 2,
\]
with the analogous identity for $p_G$.
Therefore
\[
V(D^*,G)
=
-2\log 2
+
\mathbb E_{p_{\mathrm{data}}}\left[\log \frac{p_{\mathrm{data}}}{m}\right]
+
\mathbb E_{p_G}\left[\log \frac{p_G}{m}\right].
\]
The last two terms are precisely
\[
\mathrm{KL}(p_{\mathrm{data}}\,\|\,m) + \mathrm{KL}(p_G\,\|\,m)
= 2\,\mathrm{JS}(p_{\mathrm{data}}\,\|\,p_G).
\]
Since $-2\log 2=-\log 4$, the identity follows.
\end{proof}

This is the classical divergence interpretation of GANs.
Under idealized assumptions---infinite discriminator capacity, exact optimization, and well-defined densities---training the generator is equivalent to minimizing a Jensen--Shannon divergence between the generated distribution and the data distribution.

\begin{corollary}
The global minimum of the ideal GAN objective is achieved when $p_G=p_{\mathrm{data}}$.
At that point,
\[
\mathrm{JS}(p_{\mathrm{data}}\,\|\,p_G)=0
\quad \text{and} \quad
V(D^*,G)=-\log 4.
\]
Moreover, the optimal discriminator becomes the constant function $D^*(x)=\tfrac12$.
\end{corollary}

\begin{proof}
The Jensen--Shannon divergence is nonnegative and vanishes exactly when its two arguments agree almost everywhere.
Therefore the smallest achievable value of $V(D^*,G)$ is $-\log 4$, attained when $p_G=p_{\mathrm{data}}$.
Substituting this equality into the formula for $D^*$ gives $D^*(x)=\tfrac12$.
\end{proof}

\subsection*{What is elegant in theory, and what is limited in practice}

The previous derivation is mathematically clean, but it rests on idealizations.
In practical deep learning, neither player is optimized exactly, both are represented by finite networks, and gradients are computed from mini-batches rather than population expectations.
So the Jensen--Shannon theorem is best read as a guiding idealization, not as a full theory of training dynamics.

There is also a subtle geometric issue.
If the supports of $p_{\mathrm{data}}$ and $p_G$ lie on thin and nearly disjoint manifolds, then a strong discriminator can separate them almost perfectly.
In that regime, the objective may saturate and the generator can receive weak or unstable gradients.
Thus the divergence interpretation is correct at the level of the classical objective, but it does not by itself guarantee smooth optimization.

\subsection*{A worked one-dimensional example}

Consider the toy setting in which
\[
p_{\mathrm{data}} = \tfrac12 \delta_{-1} + \tfrac12 \delta_{1},
\qquad
p_G = \delta_{\theta},
\]
so the true data has two modes at $-1$ and $1$, while the generator currently puts all mass at a single point $\theta$.
If $\theta=1$, then the generator matches only one of the two modes.
The optimal discriminator places high confidence on the unmatched real point $-1$ and reduced confidence on the shared point $1$.
The generator therefore receives a signal that it is missing part of the data distribution.

This toy example illustrates two lessons.
First, the GAN objective compares whole distributions rather than individual reconstructions.
Second, if the generator can improve its objective by concentrating on a few high-payoff regions, it may collapse onto only part of the target distribution.
This is the beginning of mode collapse.

\subsection*{Instability mechanisms}

The promise of GANs is sample quality; the cost is training instability.
Several distinct mechanisms contribute.

\paragraph{1. Saturating gradients.}
If the discriminator becomes nearly perfect, then $D_\phi(G_\theta(z))\approx 0$ for generated samples.
In the original minimax form, the generator update uses
\[
\nabla_\theta \log(1-D_\phi(G_\theta(z))),
\]
which may become very small or poorly conditioned.
A common practical modification is the non-saturating generator loss
\[
\min_\theta \; -\mathbb E_{z\sim p(z)}[\log D_\phi(G_\theta(z))],
\]
which has the same fixed point but usually stronger gradients early in training.

\paragraph{2. Mode collapse.}
The generator may map many latent values to similar outputs.
This can locally fool the discriminator while ignoring substantial parts of the data distribution.
Mode collapse is not fully explained by the ideal Jensen--Shannon derivation; it is a property of the finite, alternating optimization process.

\paragraph{3. Non-stationary targets.}
In supervised learning, the target is fixed.
In GANs, each player's target changes as the other player updates.
So optimization is inherently non-stationary.
The generator does not descend a fixed loss surface; it responds to a moving critic.

\paragraph{4. Game dynamics rather than plain minimization.}
Even simple minimax systems can rotate or cycle rather than converge directly.
This phenomenon is already visible in bilinear games such as
\[
\min_x \max_y \; xy.
\]
The point is not that GANs reduce to such a toy game, but that adversarial optimization has dynamical features absent from ordinary gradient descent.

\paragraph{5. Finite-capacity discriminators.}
The clean theorem assumes the discriminator reaches its population optimum for a fixed generator.
In practice, the discriminator is only partially optimized and represented by a finite network.
Thus the generator is updated with respect to an approximate and evolving divergence surrogate.

\subsection*{Conditional adversarial modeling}

Many applications require generation conditional on side information $c$, such as a class label, text description, segmentation map, or another image.
A conditional GAN therefore uses
\[
G_\theta(z,c), \qquad D_\phi(x,c).
\]
The classical conditional objective is
\[
\min_\theta \max_\phi
\mathbb E_{(x,c)\sim p_{\mathrm{data}}}[\log D_\phi(x,c)]
+
\mathbb E_{z\sim p(z),\,c\sim p(c)}[\log(1-D_\phi(G_\theta(z,c),c))].
\]
More carefully, one can think of the generator as trying to match the conditional data distribution $p_{\mathrm{data}}(x\mid c)$ for each context $c$.
So conditional adversarial learning is not merely adding labels; it changes the modeling target from one unconditional distribution to a family of conditional distributions.

This is why conditional GANs became important in image-to-image translation and controllable generation.
They separate two questions:
what should be generated, and under what condition should it be generated?

\subsection*{The Wasserstein viewpoint}

A major critique of the classical GAN objective is that Jensen--Shannon geometry can become poorly behaved when the model and data distributions are far apart or lie on low-dimensional sets.
The Wasserstein GAN (WGAN) addresses this by replacing the discriminator with a critic whose objective approximates an optimal-transport distance.
At a high level, one writes
\[
\max_{f\in \mathcal F_{\mathrm{Lip}}}
\left[
\mathbb E_{x\sim p_{\mathrm{data}}}[f(x)] - \mathbb E_{x\sim p_G}[f(x)]
\right],
\]
where $\mathcal F_{\mathrm{Lip}}$ denotes a class of 1-Lipschitz functions.
The generator then tries to minimize this quantity.

The significance of this reformulation is not merely technical.
It changes the geometry of distribution comparison.
The Wasserstein distance can vary continuously even when two distributions have disjoint supports, so the generator often receives more informative gradients.
This does not remove all optimization difficulty, but it explains why Wasserstein ideas became influential in adversarial learning.

\subsection*{Modeling viewpoint: explicit, implicit, and adversarial}

GANs should be understood as a special case of implicit generative modeling.
They provide a learned sampler
\[
z\sim p(z), \qquad x=G_\theta(z),
\]
without necessarily providing a tractable formula for $p_G(x)$.
Their strength is therefore sample generation rather than likelihood evaluation.
Compared with VAEs, GANs often produce sharper perceptual samples but offer less direct probabilistic structure.
Compared with explicit density models, they may avoid awkward likelihood assumptions but inherit an adversarial optimization problem.

\subsection*{What to remember}

\begin{itemize}
    \item A GAN trains a generator and discriminator through a minimax game.
    \item For a fixed generator, the optimal discriminator is
    \[
    D^*(x)=\frac{p_{\mathrm{data}}(x)}{p_{\mathrm{data}}(x)+p_G(x)}.
    \]
    \item Under ideal assumptions, substituting $D^*$ turns the generator objective into Jensen--Shannon divergence minimization.
    \item Conditional GANs target conditional distributions, while Wasserstein GANs change the geometry of distribution comparison.
    \item Practical GAN training is unstable because of saturating gradients, mode collapse, non-stationarity, and game dynamics.
\end{itemize}

\subsection*{Common confusion}

\begin{itemize}
    \item GANs do not generally provide a tractable likelihood, even though they define a model distribution implicitly.
    \item The Jensen--Shannon theorem is an idealized population statement; it is not a guarantee of easy optimization.
    \item Mode collapse is not the same as overfitting, though both can reduce diversity.
    \item Conditional GANs are not a different family of models in principle; they are adversarial models for conditional distributions.
\end{itemize}

\subsection*{Exercises}

\begin{exercise}
Derive the optimal discriminator formula by maximizing the integrand pointwise for fixed $p_{\mathrm{data}}$ and $p_G$.
\end{exercise}

\begin{exercise}
Show that if $p_G=p_{\mathrm{data}}$, then $D^*(x)=\tfrac12$ and the optimal value is $-\log 4$.
\end{exercise}

\begin{exercise}
Why can the original minimax generator loss produce weak gradients when the discriminator is very strong? Explain without computation and then verify by differentiating the loss formally.
\end{exercise}

\begin{exercise}
Compare a VAE and a GAN in terms of: (i) tractable likelihood, (ii) sampling path, (iii) training stability, and (iv) latent-space interpretation.
\end{exercise}

\begin{exercise}
In a conditional GAN, explain carefully what distribution is being matched when the condition variable $c$ is fixed.
\end{exercise}

\subsection*{Notes and references}

The original GAN framework was introduced by Goodfellow et al.
The optimal-discriminator and Jensen--Shannon derivation come from that classical paper and remain the standard theoretical entry point.
Conditional GAN ideas were developed rapidly in the years that followed and became important for controllable generation and image-to-image translation.
The Wasserstein viewpoint was introduced to address shortcomings of the Jensen--Shannon geometry in practical optimization.
For a modern textbook treatment, GANs are best read alongside latent-variable and diffusion models, because their strengths and weaknesses become clearer in comparison.
